\section{Risk-Capability Identification}
\label{sec:risk}

The discounted objective enables the actor to reason about future
$\volL$-contraction. This section develops the machinery for identifying
which capability combinations create contraction risk and when preemptive
restriction is value-positive.

\subsection{Exercise Pathways}

\begin{definition}[Exercise Pathway]
\label{def:exercise_pathway}
An \emph{exercise pathway} for a capability tuple $(d_1, \ldots, d_m)$ held
    by agent $a_k$ is a sequence of actions $(\pi_1, \ldots, \pi_L)$
    available to $a_k$ that, if executed, would produce a $\volL$-contraction
    $\Delta \volL < 0$ affecting agents beyond $a_k$. The pathway has:
\begin{itemize}
\item A \emph{prerequisite set}: the capabilities $\{d_1, \ldots, d_m\}$
    that $a_k$ must possess for the pathway to be available.
\item A \emph{contraction magnitude}: the expected $|\Delta \volL|$ if the
    pathway is fully executed.
\item An \emph{exercise probability}: $p(\Pathway \mid a_k, G_t,
    \WorldModel)$, estimated from the actor's world model $\WorldModel$.
\end{itemize}
\end{definition}

Exercise pathways are the causal chains from capability possession to
capability destruction. ``Agent $a_k$ has nuclear physics, enrichment, and
machining capabilities'' is a state. ``Agent $a_k$ builds and detonates a
weapon'' is an exercise pathway from that state, with a contraction magnitude
proportional to the number of agents destroyed and a probability conditioned
on $a_k$'s trust factor $T_k$ and behavioral history.

Not every capability combination has a dangerous exercise pathway. Most
combinations are benign: the pathways from ``can cook'' + ``can drive'' lead
only to $\volL$-neutral or expanding outcomes (meals delivered, mobility
preserved). The risk lies in combinations whose exercise pathways include
branches with large negative $\Delta \volL$. The ``beyond $a_k$''
qualifier restricts exercise pathways to externalities: self-destructive
pathways are excluded. An agent that catastrophically destroys its own
capabilities without affecting others is not captured by $\Risk$, even if
its capabilities were high-leverage or cooperative-critical. Extending
Definition~\ref{def:exercise_pathway} to include self-harm is
straightforward but deferred.

\subsection{Risk of a Capability Distribution}

\begin{definition}[Concentration Risk]
\label{def:concentration_risk}
The \emph{concentration risk} of a capability tuple $(d_1, \ldots, d_m)$
    held by agent $a_k$ at time $t$ is the expected future $\volL$-contraction
    from all exercise pathways available to $a_k$:
\begin{equation}
\label{eq:concentration_risk}
\Risk(d_1, \ldots, d_m; a_k, t) = \sum_{\Pathway}
    p(\Pathway \mid a_k, G_t, \WorldModel) \cdot |\Delta \volL(\Pathway)|
    \cdot \mathbb{E}[\gamma^{T_\Pathway}]
\end{equation}
where the sum is over all exercise pathways for the tuple, $T_\Pathway$ is
    the stochastic time to execution, and $\mathbb{E}[\gamma^{T_\Pathway}]$
    is the expected discount factor (not $\gamma^{\mathbb{E}[T_\Pathway]}$,
    which would underweight tail risk from heavy-tailed execution-time
    distributions).
\end{definition}

Concentration risk is zero for benign capability combinations (no dangerous
pathways exist) and large for combinations with high-probability, high-magnitude
exercise pathways in agents with low trust. The discount factor $\gamma$
ensures that distant risks are weighted less than imminent ones, preventing
the actor from being paralyzed by arbitrarily unlikely catastrophes in the
far future.

\paragraph{Risk depends on distribution, not possession.} Concentration of
capabilities in a single agent is itself a risk factor, related to the
power-seeking dynamics analyzed by \citet{turner2021optimal}. The same
capability tuple distributed across multiple agents who must cooperate to
exercise it has lower concentration risk than the same tuple in a single
agent, because cooperative exercise requires coordination that is observable
through the dual-channel mechanism of \citet{lasser2026poset}. Formally:
if executing pathway $\Pathway$ requires agents $a_i, a_j, a_k$ to
coordinate, the exercise probability $p(\Pathway)$ includes a factor
for coordination feasibility that is bounded by the minimum trust among
the participants and the observability of their coordination.

\subsection{When Preemptive Restriction is Value-Positive}

\begin{proposition}[Preemptive Restriction Criterion]
\label{prop:preemptive}
Under the discounted objective $\Vdisc$, restricting agent $a_k$'s access to
    a keystone capability $d_j$ in a dangerous tuple is value-positive when
    the expected future contraction from pathways requiring $d_j$ exceeds the
    immediate restriction cost:
\begin{equation}
\label{eq:preemptive}
\Risk_j(d_1, \ldots, d_m; a_k, t) > |\Delta \volL(G_t \mid \text{restrict } d_j)|
\end{equation}
where $\Risk_j$ denotes the concentration risk summed only over exercise
    pathways whose prerequisite set includes $d_j$:
\begin{equation}
\label{eq:risk_j}
\Risk_j(S; a_k, t) = \sum_{\Pathway : d_j \in \mathrm{prereq}(\Pathway)}
    p(\Pathway \mid a_k, G_t, \WorldModel) \cdot |\Delta \volL(\Pathway)|
    \cdot \mathbb{E}[\gamma^{T_\Pathway}]
\end{equation}
and the right side is the immediate $\volL$-contraction from restricting
    $d_j$ (Lemma~2 of \citet{lasser2026gfm}).
\end{proposition}

\begin{proof}
Under $\Vdisc$, the actor compares two action sequences: (1) do nothing
and face the full exercise risk; (2) restrict $d_j$ now, producing immediate
contraction $-|\Delta \volL|$ but eliminating all exercise pathways that
require $d_j$ as a prerequisite. The residual risk from pathways that do not
require $d_j$ is present under both sequences and cancels, under the
assumption that restricting $d_j$ does not change the probability of
exercise through other keystones. In practice, restriction is an observable
action that may alter $a_k$'s behavioral disposition, changing residual
pathway probabilities. The clean cancellation is an approximation. The value of
restriction exceeds the value of inaction when
$-|\Delta \volL| > -\Risk_j$, i.e., when $\Risk_j > |\Delta \volL|$.
\end{proof}

\paragraph{Connection to Lemma 2 and the precautionary principle.} The
foundational paper's Lemma~2 says restriction is always
$\volL$-contracting in isolation. The discounted objective adds the temporal
dimension: the immediate contraction is weighed against the avoided future
contraction. The preemptive restriction criterion is the quantitative
version of the quarantine pattern from Proposition~1(b) of
\citet{lasser2026gfm}, now with an explicit threshold derived from the risk
model. It also constitutes a domain-specific formalization of the
precautionary principle \citep{taleb2014precautionary}: restrict when the
expected catastrophic loss exceeds the restriction cost. Unlike the general
precautionary principle, which \citet{sunstein2005laws} argues can be
paralyzing when applied without quantitative bounds, the GFM version
provides a computable threshold
(Equation~\eqref{eq:preemptive}) grounded in the same measure that the
actor optimizes.

\paragraph{Partial pathway coverage.} If the dangerous capability tuple has
exercise pathways that do not require $d_j$ (i.e., alternate routes through
other keystone capabilities), restricting $d_j$ does not eliminate those
pathways. $\Risk_j$ captures only the risk attributable to pathways through
$d_j$. Full risk elimination may require restricting multiple keystones,
each evaluated by its own instance of
Equation~\eqref{eq:preemptive}.

\paragraph{What the actor needs.} To evaluate
Equation~\eqref{eq:preemptive}, the actor needs: (1) a world model
$\WorldModel$ that can identify exercise pathways and estimate their
probabilities; (2) trust factors $T_k$ for the agents holding dangerous
capability tuples; and (3) a contraction-magnitude estimate for the
exercise pathway. The world model is the most demanding requirement: it must
support causal counterfactual reasoning \citep{pearl2009causality} (``what
would happen if $a_k$ exercised this pathway?''), not just correlation
tracking. The companion
paper's leverage estimator (Equation~12 of \citet{lasser2026poset}) provides
a basis---leverage is already a counterfactual quantity---but extending it to
full exercise-pathway simulation is a significant computational step.
